\chapter{Representation Value: How Much Is There to Route?}
\label{ch:repr}

\chapterhandoff{Chapter~\ref{ch:system} fixed a boundary in which representation
is the only thing that varies. This chapter asks whether that variation is
worth anything. Do the modes succeed on different tasks? How much would a
chooser be worth that was told the answers afterwards? And --- the question
that turns out to decide the matter --- how much of that is worth more than
simply running one mode twice?}

\section{No mode dominates}
\label{sec:sr-per-mode}

Table~\ref{tab:sr-per-mode} gives task success for all six modes across the six
VisualWebArena cells. Read as a grid, its most useful property is a negative
one: the column that wins is not the same column everywhere.

\begin{table}[t]
\centering\small
\caption[Task success rate per mode and cell]{\textbf{Task success rate (\%) by
mode, with successes/scored-set in parentheses.} The best mode per row is in
bold. \mSoM{} wins four of six rows, but \mVis{} ties it on \cell{cls$\cdot$B2}
and \mDOM{} wins \cell{red$\cdot$B2} outright. On the WebArena cells
(Table~\ref{tab:oracle}) the winner is \mPT{} and \mDOM{}. Ordering is
cell-specific, which is the precondition for routing to be worth anything.
Cell codes name a site and a model: \cell{cls} and \cell{red} are the
classified-ads and forum sites, \cell{B0} the 235B model, \cell{B1} and
\cell{B2} the two 4B models, one from each design family.}
\label{tab:sr-per-mode}
\begin{tabular}{@{}llrrrrrr@{}}
\toprule
Cell & $n$ & \mDOM{} & \mPP{} & \mPT{} & \mPS{} & \mSoM{} & \mVis{} \\
\midrule
\cell{cls$\cdot$B0} & 224 & 17.41 & 19.64 & 15.62 & 15.62 & \textbf{27.23} & 25.00 \\
\cell{red$\cdot$B0} & 203 & 14.29 & 12.32 & 13.30 & 10.84 & \textbf{14.78} & 7.39 \\
\cell{cls$\cdot$B1} & 224 & 6.25 & 6.70 & 7.59 & 6.70 & \textbf{14.29} & 12.50 \\
\cell{red$\cdot$B1} & 203 & 5.91 & 5.42 & 5.91 & 5.91 & \textbf{7.39} & 2.46 \\
\cell{cls$\cdot$B2} & 224 & 1.34 & 1.79 & 0.45 & 0.89 & \textbf{2.23} & \textbf{2.23} \\
\cell{red$\cdot$B2} & 203 & \textbf{3.94} & 0.00 & 1.97 & 0.49 & 0.99 & 1.97 \\
\bottomrule
\end{tabular}
\end{table}

Three things follow directly from it, and each constrains what can be claimed later.

\paragraph{Base success is low, and this is the regime the dissertation is
about.} Single-mode success spans 0.45\% to 27.23\% on VisualWebArena and
reaches 35.58\% on WebArena. Every subsequent conclusion is scoped to that
range; Section~\ref{sec:regime-boundary} states what would change outside it.

\paragraph{The image-free modes do not replace the image-bearing one.} On four
of six VisualWebArena cells \mSoM{} is the strongest single mode, and on
\cell{cls$\cdot$B1} it beats the best image-free mode by 6.70\pp. Any reading of
the phantom modes as ``\mSoM{} without the cost'' is contradicted by this table,
and we do not advance one.

\paragraph{The \cell{red$\cdot$B2} row is anomalous and we do not smooth it.}
\mPP{} scores exactly zero and \mDOM{} outperforms \mSoM{} by 2.96\pp. At $n=203$
with 8 total successes, the row is dominated by counting noise, and its appearance in later analyses, including as the one cell whose learned router passes a permutation test in Chapter~\ref{ch:pred}, should be read with that
in mind.

\begin{figure}[t]
  \centering
  \includegraphics[width=\textwidth]{fig_f6_sr_heatmap.pdf}
  \caption[Success rate across modes and cells]{%
  \textbf{Success rate by mode and cell.} The same data as
  Table~\ref{tab:sr-per-mode} plus the two WebArena cells, arranged so the
  absence of a dominant column is visible at a glance. Each column is normalised
  within itself, because absolute success varies by a factor of roughly eighty
  across cells and an unnormalised map would show only which backbone is
  stronger; absolute values are printed so the transform hides nothing.}
  \label{fig:f6}
\end{figure}

\section{Cost does not track success}
\label{sec:cost-frontier}

If the expensive modes bought success in proportion to their price there would be nothing to optimise, since one would simply pay for the best affordable mode.
Figure~\ref{fig:f7} shows they do not. Within \cell{cls$\cdot$B0}, \mVis{} is
the cheapest mode at \$0.0648/episode and reaches 25.00\%, while \mPS{} costs
11.2\% more and reaches 15.62\%. Several modes are strictly dominated: they cost
more than an alternative and succeed less.

\begin{figure}[t]
  \centering
  \includegraphics[width=\textwidth]{fig_f7_cost_sr_frontier.pdf}
  \caption[Cost--success frontier per cell]{%
  \textbf{Cost against success, one panel per cell.} Non-dominated modes are
  connected; dominated modes are shown rather than hidden. Cost is
  \texttt{total\_billed\_cost\_usd} per episode for B0 and the electricity
  \texttt{total\_billed\_cost\_usd} per episode, and panels are never compared
  across the API/local boundary (Section~\ref{sec:metrics}). Seven of eight
  cells carry two or three non-dominated modes, which is the premise routing
  depends on. The exception is instructive and is discussed below.}
  \label{fig:f7}
\end{figure}

\paragraph{One cell has a single non-dominated mode, and it is the cell that
matters most later.} On \cell{cls$\cdot$B2}, \mVis{} ties \mSoM{} on success
(2.23\%) while costing 22.1\% less, and beats every other mode on both axes. The
frontier collapses to one point. That is what a regime looks like when
always-cheapest is not merely competitive but dominant, and
Chapter~\ref{ch:e2e} shows it is the general shape of the problem rather than a
quirk of one cell.

\paragraph{The price ordering is not stable across cells either.} \mVis{} is
cheapest in seven of eight cells but not on \cell{wa\_red$\cdot$B0}, where
\mDOM{} is cheapest and \mVis{} is the fourth most expensive of six. And the
spread between cheapest and dearest ranges from 11.6\% (\cell{cls$\cdot$B0}) to
77.8\% (\cell{wa\_red$\cdot$B1}), so ``the cheap mode'' and ``the expensive
mode'' name different amounts of money in different cells. Cost is an outcome of
the interaction, not a property of the representation, which is why every cost figure here is measured per episode rather than assumed from a
context-length calculation.

\section{The modes solve different tasks}
\label{sec:complementarity}

Success rates are marginals; routing needs the joint. Across
\cell{cls$\cdot$B0}, 97 of 224 tasks are solved by at least one mode, and 68 of those by more than one, so 29 tasks are solved by exactly one mode and lost
entirely if that mode is not in the portfolio. The equivalent counts for the
other cells appear in Table~\ref{tab:oracle}; the share of tasks with more than
one solver ranges from 1.5\% (\cell{red$\cdot$B2}) to 34.6\%
(\cell{wa\_red$\cdot$B0}).

The proportion of tasks that \emph{nothing} solves is the more consequential
number, and it is large: 48.1\% on the most capable cell and 92.9\% on the
weakest. Most of the benchmark is out of reach of every mode, and any policy's
room for manoeuvre is confined to the remainder.

\section{The retrospective ceiling}
\label{sec:ceiling}

Selecting per task with hindsight (Equation~\ref{eq:oracle}) gives the most that
perfect selection could ever be worth.

\begin{table}[t]
\centering\small
\caption[Oracle ceiling across all eight cells]{\textbf{The ceiling, in every
cell, on both benchmarks.} Both $\Delta$ columns compare the same thing --- an
idealised chooser told after the fact which modes solved each task --- against
the best single mode in that cell: $\Delta$SR is how much more often it
succeeds, $\Delta$cost how much less it spends when ties are broken by
measured cost. The last two columns describe the tasks themselves: the share
solved by more than one mode, and the share solved by none. The direction of
both $\Delta$ columns is the same across three models, two model families, two
benchmarks and both deployment classes.}
\label{tab:oracle}
\begin{tabular}{@{}lrlrrrr@{}}
\toprule
Cell & $n$ & best single & $\Delta$SR & $\Delta$cost & $>1$ solver & no solver \\
\midrule
\cell{Classifieds$\cdot$235B}     & 224 & \mSoM{} 27.23\% & $+16.07$\pp & $-20.2$\% & 30.4\% & 56.7\% \\
\cell{Reddit$\cdot$235B}     & 203 & \mSoM{} 14.78\% & $+11.33$\pp & $-13.7$\% & 17.7\% & 73.9\% \\
\cell{Classifieds$\cdot$Qwen-4B}     & 224 & \mSoM{} 14.29\% & $+10.27$\pp & $-30.8$\% & 12.9\% & 75.4\% \\
\cell{Reddit$\cdot$Qwen-4B}     & 203 & \mSoM{} \phantom{0}7.39\% & $+4.43$\pp & $-35.3$\% & \phantom{0}8.4\% & 88.2\% \\
\cell{Classifieds$\cdot$Gemma-4B}     & 224 & \mSoM{} \phantom{0}2.23\% & $+4.91$\pp & $-23.4$\% & \phantom{0}1.8\% & 92.9\% \\
\cell{Reddit$\cdot$Gemma-4B}     & 203 & \mDOM{} \phantom{0}3.94\% & $+3.45$\pp & $-26.6$\% & \phantom{0}1.5\% & 92.6\% \\
\midrule
\cell{Reddit (WA)$\cdot$235B} & 104 & \mPT{} 35.58\% & $+16.35$\pp & $-26.9$\% & 34.6\% & 48.1\% \\
\cell{Reddit (WA)$\cdot$Qwen-4B} & 104 & \mDOM{} 16.35\% & $+14.42$\pp & $-33.5$\% & 17.3\% & 69.2\% \\
\bottomrule
\end{tabular}
\end{table}

\begin{scopebox}
\textbf{Headline, with its comparator and its cell set attached.} A
retrospective oracle over six observation modes improves task success by
$+3.45$ to $+16.35\pp$ over the best single mode while spending 1.6--35.3\%
less, in \textbf{8 of 8} cells spanning two benchmarks. The upper bound
$+16.35\pp$ is \cell{wa\_red$\cdot$B0}; restricted to the six VisualWebArena
cells the upper bound is $+16.07\pp$ (\cell{cls$\cdot$B0}). Both figures are
correct and they index different cell sets, so neither may be quoted without
saying which.

The cost range needs its own note, because it moved during this work and the
reason is instructive. ``Best single mode'' is ambiguous when two fixed modes tie
on success, and two cells do: \cell{cls$\cdot$B2} (\mSoM{} and \mVis{} both at
2.23\%) and \cell{wa\_red$\cdot$B1} (\mDOM{}, \mPT{} and \mPP{} all at
16.35\%). Breaking those ties by the mode list's order selects the
\emph{dearer} tied mode, which flatters the oracle's cost advantage: on
\cell{cls$\cdot$B2} it reports $-23.4$\% where a cost-aware tie-break gives
$-1.6$\%. We report the cost-aware figures throughout. \textbf{The range is
therefore 1.6--35.3\%, not the 13.7--35.3\% a list-order comparator would
produce.} The success side is unaffected, since tied modes have equal success by definition, and the direction still holds in 8 of 8 cells.

\textbf{This number travels with two qualifications for the rest of the dissertation, and they are not cosmetic.} Adding a single \emph{rerun} arm of an
unchanged mode buys 2.0--7.6\pp{} by itself
(Section~\ref{sec:noise-floor}). And at matched arm count, six reruns of the
strongest single mode would reach within 3.53\pp{} of the six-mode oracle, below our own noise threshold (Section~\ref{sec:six-vs-six}). What survives is
not that representation selection reaches a \emph{higher} ceiling than
repetition, but that it reaches a comparable one at roughly a sixth of the
serving cost.
\end{scopebox}

\subsection{Where the ceiling comes from}

The oracle decomposes into the two halves defined in
Section~\ref{sec:router-formulation}, and they behave very differently.

\emph{Triage}, which keeps the best-SR mode but spends the cheapest on tasks that nothing will solve, reduces cost by up to 30.6\% at exactly zero success
change. The zero is by construction, not a discovery: the policy only redirects
tasks that no mode solves, so no success can be lost.

\begin{scopebox}
\textbf{The same tie-break correction applies here, and it costs one cell.}
Measured against the list-order ``best single mode'', triage saves 9.5--30.6\%
in 8 of 8 cells. Against the \emph{cost-aware} tied best single
(Section~\ref{sec:ceiling}) the range is $-30.6$\% to $+1.1$\%, cheaper in
\textbf{7 of 8}. The cell that flips is \cell{cls$\cdot$B2}, where \mVis{}
matches \mSoM{}'s 2.23\% success for 22.1\% less money: against that comparator
triage is 1.1\% \emph{more} expensive, so its apparent saving there was an
artefact of comparing to a needlessly expensive fixed policy. We quote the
corrected figures. The direction survives in 7 of 8 and the upper end is unchanged, but ``in every cell'' does not survive, and it was the strongest
word in the sentence.
\end{scopebox}
Its interest lies elsewhere: its label is binary, defined on every task, and
therefore \emph{never supply-constrained}. Chapter~\ref{ch:pred} tests exactly
this half, which is the most favourable version of the problem we can pose.

\emph{Which-mode}, choosing among the modes that do solve a task, supplies the entire success gain. Its label exists only on solved tasks: 97 of 224 on
\cell{cls$\cdot$B0}, and 15 on \cell{red$\cdot$B2}. That asymmetry is the seed
of Chapter~\ref{ch:disc}'s mechanism.

The full ladder of policies per cell --- every fixed mode, both halves of the
oracle, the version that breaks ties by measured cost, and the policy that
tries modes in increasing order of cost until one succeeds --- is given in
Appendix~\ref{app:ceiling}, which also isolates how much of the oracle's cost
advantage comes from the tie-break rule alone.

\subsection{What the ceiling is not}

It is not a policy. For each task it reads off the outcome of every mode,
including modes that were never run on that task --- information no deployed
system can have.

It is also not free of its own label noise, and the amount is measurable rather
than rhetorical. Section~\ref{sec:noise-floor} measures 12.05--14.29\%
task-level discordance between two runs of an \emph{identical} condition, from
which Section~\ref{sec:six-vs-six} estimates that roughly 24\% of that cell's
tasks are \emph{flippable} for a single mode: they can come out either way on
two runs of the same setup, with nothing changed between them. The oracle's
``this mode solves this task'' labels inherit that instability, so a live rerun
would not reproduce its choices task for task. Its \emph{triage} result is where
this bites hardest: the ``zero success loss'' is zero on \emph{this} log only,
because the tasks it sends to the cheap mode are the ones that happened to
fail in this run, not the ones that would fail on average.

\begin{figure}[t]
  \centering
  \includegraphics[width=\textwidth]{fig_f8_oracle_ceiling.pdf}
  \caption[The ceiling with its repetition baseline]{%
  \textbf{Oracle gain per cell, with the measured rerun band drawn alongside.}
  Bars are the six-mode oracle gain over the best single mode --- that is,
  what \emph{five added arms} buy when each added arm is a different mode
  rather than a repeat. The shaded band is what \emph{one} rerun arm of an
  unchanged mode buys, measured on replicated conditions
  (Section~\ref{sec:noise-floor}). The band is drawn beside the bars and not
  subtracted from them: the two quantities differ in arm count, and the
  subtraction would be exactly the arithmetic our own noise inventory forbids.
  The like-for-like one-arm comparison is Figure~\ref{fig:f10b}.}
  \label{fig:f8}
\end{figure}

\section{Each mode's irreplaceable contribution}
\label{sec:dropone}

Drop-one (Equation~\ref{eq:dropone}) asks what removing a mode from the
portfolio costs. Table~\ref{tab:dropone} gives all 36 values with bootstrap
intervals.

\begin{table}[t]
\centering\small
\caption[Drop-one contribution per mode and cell]{\textbf{Drop-one loss
(\pp) with bootstrap 95\% intervals.} Positive means the mode solves something
no other mode solves. Direction is consistent, with 32 of 36 values positive and 14 of 18 among the phantom modes, but magnitude is small: only \emph{one}
value clears the exact one-sided 95\% cutoff for its own arm
(Section~\ref{sec:floor-spread}), and it is a baseline mode; three clear the looser normal approximation. \emph{All 36 rows are} \texttt{complete\_exact}
\emph{and non-partial; the} \texttt{NON\_PAPER\_GRADE} \emph{stamp in the
artefact header records the producer's invocation mode, not a defect in these
rows, and is reported here rather than quietly dropped or quietly relied upon.}
$^\dagger$\emph{Four cells observe zero tasks solved by that mode alone. The
usual resampling interval breaks down there: every resample of an all-zero
column is also zero, so it returns} $[0,0]$\emph{, a certainty the data cannot
supply. Those four upper
limits are exact two-sided binomial bounds instead: 1.63\pp{} at} $n=224$
\emph{and 1.80\pp{} at} $n=203$.}
\label{tab:dropone}
\begin{tabular}{@{}lcccccc@{}}
\toprule
Cell & \mDOM{} & \mPP{} & \mPT{} & \mPS{} & \mSoM{} & \mVis{} \\
\midrule
\cell{Classifieds$\cdot$235B} & 1.79 & 2.68 & 0.89 & 0.89 & 2.68 & \textbf{4.02} \\
                    & \tiny[0.00,\,3.57] & \tiny[0.89,\,4.91] & \tiny[0.00,\,2.23] & \tiny[0.00,\,2.23] & \tiny[0.89,\,4.91] & \tiny[1.79,\,6.70] \\
\cell{Reddit$\cdot$235B} & 1.48 & 0.99 & 0.99 & 0.99 & 1.97 & 1.97 \\
                    & \tiny[0.00,\,2.96] & \tiny[0.00,\,2.46] & \tiny[0.00,\,2.46] & \tiny[0.00,\,2.46] & \tiny[0.49,\,3.95] & \tiny[0.49,\,3.94] \\
\cell{Classifieds$\cdot$Qwen-4B} & 0.45 & 0.89 & 0.45 & 1.34 & \textbf{4.46} & \textbf{4.02} \\
                    & \tiny[0.00,\,1.34] & \tiny[0.00,\,2.23] & \tiny[0.00,\,1.34] & \tiny[0.00,\,3.13] & \tiny[1.79,\,7.14] & \tiny[1.77,\,7.14] \\
\cell{Reddit$\cdot$Qwen-4B} & 0.49 & 0.99 & 0.49 & 0.00 & 0.49 & 0.99 \\
                    & \tiny[0.00,\,1.48] & \tiny[0.00,\,2.46] & \tiny[0.00,\,1.48] & \tiny[0.00,\,1.80]$^\dagger$ & \tiny[0.00,\,1.48] & \tiny[0.00,\,2.46] \\
\cell{Classifieds$\cdot$Gemma-4B} & 0.45 & 0.00 & 0.00 & 0.45 & 2.23 & 2.23 \\
                    & \tiny[0.00,\,1.34] & \tiny[0.00,\,1.63]$^\dagger$ & \tiny[0.00,\,1.63]$^\dagger$ & \tiny[0.00,\,1.34] & \tiny[0.45,\,4.46] & \tiny[0.45,\,4.03] \\
\cell{Reddit$\cdot$Gemma-4B} & 2.96 & 0.00 & 0.49 & 0.49 & 0.49 & 1.48 \\
                    & \tiny[0.99,\,5.42] & \tiny[0.00,\,1.80]$^\dagger$ & \tiny[0.00,\,1.48] & \tiny[0.00,\,1.48] & \tiny[0.00,\,1.48] & \tiny[0.00,\,3.45] \\
\bottomrule
\end{tabular}
\end{table}

\begin{scopebox}
\textbf{This is a downgraded claim, and the downgrade is ours.} An earlier reading of these data, reported in our own project record, described the phantom modes as forming a structurally independent routing region, on the basis
that their drop-one contributions were consistently positive. Consistency of
\emph{direction} survives: 14 of 18 phantom arm--cell values are positive. But
the two axes these modes separate --- which page text is supplied, and how the
instructions are worded --- pool to effects of only $1.35\pm0.32$\pp{} and
$2.09\pm0.43$\pp{} (fixed-effect pool over the six cells). Those sit four to five standard errors from zero, so the structure is not a null result, and still below the most permissive repetition floor we measure (2.0\pp). Only 1 of 36 drop-one values clears the exact one-sided cutoff at which a single
rerun would be unlikely to produce the effect on its own; three do under the looser normal approximation, which is why Section~\ref{sec:floor-spread} adopts
the exact null. Precision is not the problem; magnitude is. \textbf{The correct statement is therefore: we
observe a directionally consistent complementarity structure whose magnitude
does not exceed the repeated-sampling floor.} It is not advanced as a positive
structural claim.
\end{scopebox}

Figure~\ref{fig:f9} makes the reason legible. The four text-side modes overlap
heavily with each other because they differ only in text format and prompt
style; the larger unique contributions come from \emph{crossing sides}: \mSoM{} and \mVis{}, the only modes that send an image. Grouping the
modes by what is actually transmitted turns an otherwise puzzling pattern into
an expected one.

\begin{figure}[t]
  \centering
  \includegraphics[width=\textwidth]{fig_f9_drop_one_vs_floor.pdf}
  \caption[Drop-one contribution against the noise threshold]{%
  \textbf{Every drop-one value against the threshold it must clear.} Modes are
  grouped by what is sent: four text-side modes (no image), one combined, one
  visual. The vertical band is the one-sided 95\% threshold of
  Section~\ref{sec:floor-spread}: 3.82--4.15\pp{} under the normal
  approximation the figure draws, and 4.02--4.46\pp{} under the exact binomial
  null the text adopts, which is stricter still. Almost every value sits left of
  it, and every value that clears it is a baseline mode. The figure is built so that
  consistent direction and insufficient magnitude cannot be mistaken for one
  another.}
  \label{fig:f9}
\end{figure}

\section{What running the same thing twice buys}
\label{sec:noise-floor}

The measurements above are differences between modes. To interpret them we need
to know how large a difference appears between two runs of the \emph{same}
mode. We therefore replicated identical conditions and compared their solve
sets.

\subsection{Measured discordance}

Three arms of \cell{cls$\cdot$B0} have been run twice under an unchanged
configuration, plus one WebArena pilot/full overlap
(Table~\ref{tab:noise-floor}).

\begin{table}[t]
\centering\small
\caption[Same-condition run-to-run floors]{\textbf{Two runs of one unchanged
condition.} $\mathrm{self\_drop}(a\!\to\!b) = |\{a \text{ solves}\} \setminus
\{b \text{ solves}\}| / n$, the same functional as Equation~\ref{eq:dropone}, so these are directly comparable to a one-arm
addition. Discordance is the fraction of tasks on which the two runs disagree in
either direction.}
\label{tab:noise-floor}
\begin{tabular}{@{}llrrrr@{}}
\toprule
Pair & scope & $n$ & $a\!\to\!b$ & $b\!\to\!a$ & discordance \\
\midrule
\cell{B0.cls.dom}    & B0 $\times$ classifieds & 224 & 7.14\pp & 4.91\pp & 12.05\% \\
\cell{B0.cls.vision} & B0 $\times$ classifieds & 224 & 7.59\pp & 6.70\pp & 14.29\% \\
\cell{B0.cls.som}    & B0 $\times$ classifieds & 224 & 5.36\pp & 7.59\pp & 12.95\% \\
\cell{B1.wa-red}     & 5 modes $\times$ 10 shared tasks & 50 & 2.00\pp & 4.00\pp & \phantom{0}6.00\% \\
\bottomrule
\end{tabular}
\end{table}

Two runs of an identical condition disagree on 12--14\% of tasks. The agent's
decoding is greedy, so this is not sampling stochasticity; it is environment
drift, session state and timing. We had previously recorded the belief that
greedy decoding made local backbones bit-identical across reruns. Step-level
determinism is real and was verified; \emph{episode}-level determinism does not
follow, because an episode also depends on site state, wall-clock and session lifetime. Table~\ref{tab:noise-floor} refutes the stronger belief, and it was
our own.

\subsection{One draw is not a bound}
\label{sec:floor-spread}

The measured $|\Delta\mathrm{SR}|$ between two runs of one condition spans
0.89--2.23\pp. It is tempting to call that the noise floor and declare anything
larger to be signal. That reasoning compares a draw to a draw.

Under the exchangeability null, where the two runs are the same condition and each discordant task flips either way with probability $\tfrac12$, the difference is $D=(2X-d)/n$ with $X\sim\mathrm{Binom}(d,\tfrac12)$, giving
\begin{equation}
  \mathrm{SD}(D)=\frac{\sqrt{d}}{n},
  \label{eq:floor-sd}
\end{equation}
where $d$ is the discordant count. Table~\ref{tab:floor-spread} evaluates it.

\begin{table}[t]
\centering\small
\caption[The spread of the repetition difference]{\textbf{What repetition
\emph{could} deliver, not only what it did.} The observed $|\Delta$SR$|$ column
is one draw of a random quantity whose own standard deviation is the next
column. The one-sided 95\% point is the threshold an effect must reach before a
single rerun would be unlikely to produce it alone.}
\label{tab:floor-spread}
\begin{tabular}{@{}lrrrrrr@{}}
\toprule
Pair & $n$ & discordant $d$ & obs.\ $|\Delta$SR$|$ & SD under null
  & 95\% (normal) & \textbf{95\% (exact)} \\
\midrule
\cell{B0.cls.dom}    & 224 & 27 & 2.23\pp & 2.32\pp & 3.82\pp & \textbf{4.02\pp} \\
\cell{B0.cls.vision} & 224 & 32 & 0.89\pp & 2.53\pp & 4.15\pp & \textbf{4.46\pp} \\
\cell{B0.cls.som}    & 224 & 29 & 2.23\pp & 2.40\pp & 3.95\pp & \textbf{4.02\pp} \\
\bottomrule
\end{tabular}
\end{table}

The upper edge of the observed band (2.23\pp) is the same order as one standard
deviation (2.32--2.53\pp). \textbf{Clearing the band is therefore not clearing
the noise.}

How much further an effect must reach depends on whether the null is evaluated
exactly or approximated, and the difference is not cosmetic. Multiplying
$\mathrm{SD}(D)$ by 1.645 gives 3.82--4.15\pp, but $D$ can only take a
discrete set of values, fixed by how many tasks the two runs disagree on,
rather than varying continuously; at these values of $d$ the normal
approximation is optimistic. Evaluating $X\sim\mathrm{Binom}(d,\tfrac12)$
exactly gives one-sided 95\% cutoffs of \textbf{4.02, 4.02 and 4.46\pp} for
$d=27,29,32$. We use the exact values. Of the 36 drop-one values in
Table~\ref{tab:dropone}, three clear the normal approximation and only
\textbf{one} clears the exact cutoff for its own arm.

Both figures are themselves lower bounds, since Equation~\ref{eq:floor-sd}
assumes only exchangeability and does not model the one-directional environment
drift we also observe. Appendix~\ref{app:noise} gives the derivation, separates
the two distinct floor quantities this literature tends to conflate, and states
which comparisons each one licenses.

\begin{figure}[t]
  \centering
  \includegraphics[width=0.72\textwidth]{fig_f10_rerun_discordance.pdf}
  \caption[Run-to-run discordance on identical conditions]{%
  \textbf{Two runs of the same condition disagree on 12--14\% of tasks.}
  Per-pair discordance for the three replicated \cell{cls$\cdot$B0} arms.
  Decoding is greedy; this is environment and session variation, not sampling.}
  \label{fig:f10}
\end{figure}

\subsection{The like-for-like comparison}
\label{sec:one-arm-margin}

The honest test is arm-count matched. Starting from one mode, what does adding \emph{one} more arm buy, if that arm is a different representation rather than a rerun of the same representation? Both are
$|\{\text{added}\}\setminus\{\text{baseline}\}|/n$, the same functional at the
same arm count on the same tasks.

On \cell{cls$\cdot$B0}, starting from \mSoM{} at 27.23\%: adding the best
distinct representation (\mDOM{}) buys $7.14\pp$; adding a rerun buys
$4.91$--$7.59\pp$. The representation arm is \emph{inside} the rerun band.

\begin{figure}[t]
  \centering
  \includegraphics[width=\textwidth]{fig_f10b_one_arm_margin.pdf}
  \caption[Representation arm versus rerun arm at matched arm count]{%
  \textbf{Adding one extra arm, either a different representation or the same one again.}
  Both bars are $|\{\text{added}\}\setminus\{\text{baseline}\}|/n$ on the same
  cell, same $n$, same functional, same arm count, which is the only comparison the data licenses. On \cell{cls$\cdot$B0} the distinct representation lands inside
  the rerun band. The formal relationship to \mbox{pass@$k$}
  \citep{kulal2019spoc, chen2021codex} is direct: adding a rerun arm is
  \mbox{pass@2} of one mode, and the question is whether a second
  \emph{representation} outperforms a second \emph{attempt}.}
  \label{fig:f10b}
\end{figure}

\begin{scopebox}
\textbf{What the one-arm result licenses.} \emph{Licensed} on
\cell{cls$\cdot$B0}: at the one-arm margin, a distinct representation is worth
no more than a rerun of the same representation. \emph{Not licensed} from this
comparison alone: ``the whole six-mode ceiling is noise''. The ceiling adds
\emph{five} arms; we hold \emph{one} rerun arm, and reruns have their own
diminishing returns, so the subtraction is not available. The next section asks
the six-arm question directly instead of declining it.
\end{scopebox}

\section{Six reruns versus six modes}
\label{sec:six-vs-six}

Refusing an invalid subtraction is not the same as answering the question behind
it. If one extra arm of repetition buys up to 7.6\pp, what would \emph{five} buy, and is the six-mode ceiling distinguishable from six repetitions of the
strongest single mode? We cannot measure this directly, since that would need six
replicate runs of one condition. We can bound it.

Adopt the same exchangeability null already used in
Section~\ref{sec:floor-spread}, and nothing further. Partition the scored set, for one mode, into $s_0$ tasks it solves deterministically, $m$ tasks it
solves with probability $\tfrac12$, and the rest it never solves. Two runs then
disagree on each flippable task with probability $\tfrac12$, so the observed
discordant count $d$ estimates $m/2$, giving $\hat m = 2d$ and
$\hat s_0 = |A| - d$. The union over $k$ independent runs, which is $\mathrm{pass}@k$ for that one arm, is
\begin{equation}
  U(k) = s_0 + m\left(1 - 2^{-k}\right).
  \label{eq:union}
\end{equation}

$U(1)$ reproduces the single-run success rate by construction, so it is not
evidence. $U(2)$ is not: it can be checked against the \emph{observed} two-run
union, and Table~\ref{tab:six-vs-six} reports both.

\begin{table}[t]
\centering\small
\caption[Six reruns of one mode against six distinct modes]{\textbf{What
repetition alone would have reached, at the same arm count.} $U(2)$ model versus
observed is an out-of-sample check on Equation~\ref{eq:union}; it errs by at
most $1.12$\pp{} in either direction. $U(6)$ is an \emph{upper} bound on
repetition: treating every discordant task as a fair coin maximises how fast the
union grows, and flippable tasks with $p\neq\tfrac12$ saturate sooner.
The classified-ads site with the 235B model (\cell{cls$\cdot$B0}) is the only
cell for which a mode was run twice.}
\label{tab:six-vs-six}
\begin{tabular}{@{}lrrrrr@{}}
\toprule
Replicated arm & single run & discordant & $U(2)$ model & $U(2)$ observed & $\mathbf{U(6)}$ \\
\midrule
\mDOM{}   & 17.41\% & 27 & 23.44\% & 22.32\% & \textbf{29.09\%} \\
\mVis{}   & 25.00\% & 32 & 32.14\% & 31.70\% & \textbf{38.84\%} \\
\mSoM{}   & 27.23\% & 29 & 33.71\% & 34.82\% & \textbf{39.77\%} \\
\midrule
\multicolumn{5}{@{}l}{six-\emph{mode} oracle on the same cell}
  & \textbf{43.30\%} \\
\bottomrule
\end{tabular}
\end{table}

Six reruns of \mSoM{} would reach 39.77\%; six distinct modes reach 43.30\%.
The residual is \textbf{3.53\pp}.

\begin{scopebox}
\textbf{This is the strongest attack on our headline, and it substantially
lands.} $3.53$\pp{} sits \emph{below} the threshold this dissertation derives in
Section~\ref{sec:floor-spread} and applies to everyone else's effects: below it under the normal approximation ($3.82$--$4.15$\pp), and further below it
under the exact binomial null ($4.02$--$4.46$\pp) that we adopt. Applying our own standard to our own headline: \textbf{at
matched arm count, six distinct representations are not distinguishable from six
repetitions of the strongest one.} Most of the ceiling is repetition.

\textbf{What survives is a different claim, on a different axis.} The two are not
interchangeable at deployment. The mode oracle spends \emph{one} episode per task, selecting a representation rather than retrying. The six-rerun union
spends \emph{six}. So representation selection reaches a comparable ceiling at
roughly one-sixth of the serving cost. The defensible headline is therefore
about the \emph{cost-efficiency} of the ceiling, not its height.
\end{scopebox}

The analysis has two limits, both of which cut against reading it too strongly in either direction. It rests on one cell, because \cell{cls$\cdot$B0} is the only
cell with replicated arms; and Equation~\ref{eq:union} compresses a whole
distribution of task-level failure probabilities into just two estimated
numbers. Its only defence is that we check it at $k=2$ against what actually
happened, and that it errs towards the conclusion we did not want.

\section{What the WebArena cells add, and what they do not}
\label{sec:wa-scope}

The two WebArena cells extend the ceiling result to a second benchmark and, in
\cell{wa\_red$\cdot$B0}, produce its largest value ($+16.35\pp$). They also
show the mode ordering rearranging: the best single mode there is \mPT{}
(35.58\%), an image-free mode, with \mSoM{} fourth. That is consistent with the
corpus difference in Figure~\ref{fig:f4}(B), but consistency is not attribution, and we make the weaker claim
(Section~\ref{sec:benchmarks}).

The WebArena rerun figure in Table~\ref{tab:noise-floor} must be read with more
care than the others, and we flag it rather than averaging it in. It pools
five modes' pilot-versus-full comparisons over the same 10 registered pilot
tasks, so its 50 comparisons come from only ten distinct tasks, while the distinct-arm figure
beside it is computed on 104 tasks starting from \mDOM{}, whose own pilot-versus-full comparison shows zero flips in either direction. The 2--4\pp{}
band is generated entirely by \mVis{} and \mPP{}. It should be read as
\emph{some arms of this cell move by 2--4\pp{} under repetition}, not as a floor
for \mDOM{}.

\section{What this chapter establishes}
\label{sec:ch4-summary}

Three findings, in descending order of strength.

\textbf{Strong.} A retrospective oracle over the six modes gains
$+3.45$ to $+16.35\pp$ at 1.6--35.3\% lower cost, in 8 of 8 cells across two
benchmarks, three backbones and two model families. The direction is not in
doubt.

\textbf{Qualified, and the qualification is severe.} If only one extra run is
allowed, a different representation is worth no more than a repeat of the same
one. At the matched arm count of six, repeating the strongest mode comes
within 3.53\pp{} of what six different modes achieve, which is below the
threshold this chapter derived. The two sides are matched in arms, not in what
they cost to run: the mode oracle reaches its level with one episode per task,
where repetition needs six. The height of
the ceiling is largely a repeated-sampling effect. Its \emph{cost} is not: the
mode oracle reaches that height for one episode per task where repetition needs
six. The surviving claim is about the price of the ceiling, not its height.

\textbf{Downgraded.} The three phantom modes contribute in a consistent
direction (14 of 18 arm--cell values positive) but at a magnitude below the
repetition floor. We report the structure and decline the structural claim.

\chapterhandoff{Comparing the six fixed modes establishes that there is real
room above the best single one, and that a meaningful share of that room can
be bought by repetition instead of by a different representation. But the
chooser that reaches it is told which modes worked, which no deployed system
can know. The operational question is whether the need for a different
representation shows up in anything observable \emph{before} the expensive
choice is made. Chapter~\ref{ch:pred}
formulates that prediction problem and tests it under a protocol designed so
that a negative answer cannot be blamed on test information having leaked into
training, and a positive answer cannot be produced by trying decision
thresholds until one works.}
